\section{What the Greeks Reached}

\subsection{Potency and act}

Aristotle's contribution begins with an observation about change so plain that its power is easy to miss. Whatever changes acquires something it did not have: the cold water becomes hot, the seed becomes an oak, the unlit lamp lights. Before the change, the thing possessed the new state only as a capacity---Aristotle says \emph{potentially}; after, it possesses the state in fact---\emph{actually}. And nothing confers on itself what it does not have. The water cannot heat itself; something already hot, or already in act in the relevant way, must do it. Potency does not lift itself into act, any more than an empty pocket funds a purchase.

Aristotle states the rule in a sentence that governs everything downstream, and states it about ordinary cases: ``from the potential the actual is always produced by an actual thing, e.g.\ man by man, musician by musician; there is always a first mover, and the mover already exists actually.''\endnote{Aristotle, \emph{Metaphysics} IX ($\Theta$), ch.~8 (1049b), in the public-domain translation of W.~D. Ross, \emph{The Works of Aristotle Translated into English}, vol.~VIII: \emph{Metaphysica} (Oxford: Clarendon Press, 1908); wording checked 2026-07-25 in the Internet Archive scan of the 1908 printing (item \texttt{worksofaristotle0008unse}), re-downloaded that date and found byte-identical to the copy checked 2026-07-24. The same chapter adds the illustrations Aristotle takes from craft: ``it is thought impossible to be a builder if one has built nothing, or a harper if one has never played the harp; for he who learns to play the harp learns to play it by playing it.''} The examples are deliberately unglamorous---a man, a harpist---because the claim is not yet about first causes but about causing as such. Give the principle its widest form and it says: no capacity ever cashes itself; every actualization traces to something already actual.

From this Aristotle draws his ranking, and the ranking is more careful than admirers of the phrase usually notice. Actuality is prior to potency, he says, ``both in formula and in substance; and in time it is prior in one sense, and in another not.''\endnote{\emph{Metaphysics} IX, ch.~8 (1049b), Ross, checked 2026-07-25 as above. The temporal qualification is Aristotle's own: ``the actual member of a species is prior to the potential member of the same species, though the individual is potential before it is actual''---this seed precedes this oak, but an oak preceded the seed.} In definition, because one cannot say what a capacity is without naming what it is a capacity \emph{for}: sight is explained by seeing, not seeing by sight. In time, only with the qualification he insists on---this seed precedes this oak, but an oak preceded the seed. And in substance, decisively: ``matter exists in a potential state, just because it may attain to its form,'' so that the capable exists for the sake of the achieved and never the reverse.\endnote{\emph{Metaphysics} IX, ch.~8 (1050a--b), Ross, checked 2026-07-25 as above: ``men have the art of building that they may build\ldots{} Further, matter exists in a potential state, just because it may attain to its form; and when it exists actually, then it is in its form.'' The chapter's closing inference---``actuality is prior in substantial being to potency; and as we have said, one actuality always precedes another in time right back to the actuality of the eternal prime mover''---is the sentence that carries Aristotle's cosmology.} The chapter closes by drawing the line all the way out: actuality ``is prior in substantial being to potency; and as we have said, one actuality always precedes another in time right back to the actuality of the eternal prime mover.''

\subsection{The series that must have a first}

Follow that pointing finger through a series. The water is heated by the fire, the fire sustained by its fuel and air, and so on. Where the members of such a series act \emph{simultaneously}---each moving only because something is now moving it---the series cannot do without a first member whose actuality is not on loan. Deferral is not explanation; a chain of borrowers, however long, never adds up to an owner.

Aristotle argues the point in its own right, early, before any theology is in view. ``Evidently there is a first principle,'' he writes, ``and the causes of things are neither an infinite series nor infinitely various in kind''---and he runs the case through each of his four causes in turn, refusing an endless regress in matter, in efficient causing, in ends, and in definition.\endnote{\emph{Metaphysics} II ($\alpha$), ch.~2 (994a), Ross, checked 2026-07-25 as above. Aristotle's efficient-cause example in the same passage is a chain of simultaneous movers---``man for instance being acted on by air, air by the sun''---and his verdict on the endless series of ends is unusually warm: ``those who maintain the infinite series destroy the Good without knowing it. Yet no one would try to do anything if he were not going to come to a limit.'' The borrower-and-owner image is this essay's, not Aristotle's.} His remark about the endless series of ends carries the whole temperament of the argument: ``those who maintain the infinite series destroy the Good without knowing it. Yet no one would try to do anything if he were not going to come to a limit.'' An explanation that never lands is not a long explanation. It is a promise repeatedly renewed and never kept.

In the twelfth book he applies the rule to the cosmos. Eternal substances are not enough, he notes, and neither is a substance that merely \emph{could} act: ``if it does not act, there will be no movement. Further, even if it acts, this will not be enough, if its essence is potency\ldots{} There must, then, be such a principle, whose very essence is actuality.''\endnote{\emph{Metaphysics} XII ($\Lambda$), ch.~6 (1071b), Ross, checked 2026-07-25 as above. The argument is aimed partly at the Platonists: eternal Forms explain nothing about movement, ``unless there is to be in them some principle which can cause movement.''} Then, in the following chapter, the terminus is named: ``there is a mover which moves without being moved, being eternal, substance, and actuality.'' And then, in a sentence whose reach exceeds its grammar: ``On such a principle, then, depend the heavens and the world of nature.''\endnote{\emph{Metaphysics} XII, ch.~7 (1072a--b), Ross, checked 2026-07-25 as above. The same chapter supplies the celebrated description of the first principle's life: ``And life also belongs to God; for the actuality of thought is life, and God is that actuality\ldots{} We say therefore that God is a living being, eternal, most good.'' The scholastic term for a series whose members cause only while the prior member is causing---an \emph{essentially ordered} series---is not Aristotle's; the distinction is drawn out in section~4 below.} \emph{Depend} is his word. The world of nature, for Aristotle, is not a house standing after its builder; it hangs, now, from pure actuality.

\subsection{Plato: what is had, and what is}

Plato had already reached a structurally similar result from a different side. What strikes him is not change but multiplicity: many beautiful things, no one of which is beauty. Each is beautiful in some respect, for a while, from an angle; each has beauty the way a borrower has money. In the \emph{Phaedo} Socrates offers what he calls the safest answer to the question why anything is beautiful: ``nothing makes a thing beautiful but the presence and participation of beauty,'' so that ``by beauty all beautiful things become beautiful.''\endnote{Plato, \emph{Phaedo} 100c--e, in the public-domain translation of Benjamin Jowett; wording checked 2026-07-25 at Project Gutenberg (ebook no.~1658, re-downloaded that date). Socrates prefaces the formula with its condition: ``if there be anything beautiful other than absolute beauty\ldots{} it can be beautiful only in as far as it partakes of absolute beauty.'' The passage belongs to his account of the ``second sailing''; its generalization from beauty to being is the later Platonic and Christian-Platonic development, not a claim of the dialogue itself.} Strip the example and keep the structure: whatever possesses a perfection partially, derivatively, and in company with many others that possess it too, points beyond itself to something that does not \emph{have} that perfection but simply \emph{is} it.

In the \emph{Republic} Plato pushes the same structure one step further, and the step is worth quoting for where it lands. Explaining the Good by comparison with the sun---the sun which is ``not only the author of visibility in all visible things, but of generation and nourishment and growth''---he says that ``the good may be said to be not only the author of knowledge to all things known, but of their being and essence, and yet the good is not essence, but far exceeds essence in dignity and power.''\endnote{Plato, \emph{Republic} VI, 508e--509b, Jowett translation; wording checked 2026-07-25 at Project Gutenberg (ebook no.~1497). Two cautions belong with the quotation. First, the passage is an extended analogy about the conditions of knowledge, not a doctrine of creation: the Good is not said to make anything exist from nothing, and no maker appears. Second, the clause rendered ``far exceeds essence in dignity and power'' (\latin{epekeina tēs ousias}) has been read in more than one way since antiquity, and nothing here depends on the stronger readings. It is cited for the structure it displays---source of being, itself outside the class of things that have being---and for the fact that the sun image was already in place at this exact joint of the argument.} Notice which picture is doing the work at that precise joint: the sun, the source that gives what the lit things do not own. The lamp and the lit air of the first pages are not a modern convenience. They were the Greek picture, waiting.

\subsection{What the Greeks did not reach}

It matters, for honesty and for everything that follows, to state precisely what the Greeks did \emph{not} conclude.

Aristotle's unmoved mover is an explanation of motion, not of being, and he says exactly how it moves: ``The final cause, then, produces motion by being loved, and by that which it moves, it moves all other things.''\endnote{\emph{Metaphysics} XII, ch.~7 (1072b), Ross, checked 2026-07-25 as above. The first principle moves as an object of desire and thought---as the beloved moves the lover---not as an efficient cause conferring existence. Whether Aristotle's mover is in some further sense an efficient cause of the cosmos is long disputed in the scholarship; this essay takes no position beyond the text's own statement and builds nothing on the disputed reading.} It moves the heavens as the beloved moves the lover: by being desirable, not by supplying anything. There is no suggestion that it chose the world, knows the world's particulars, or gave the world existence. The cosmos alongside it is eternal, underived, and simply there.

Plato's forms are patterns, not makers, and where a maker does appear in his writing, in the \emph{Timaeus}, the limits sit on the surface of the text. The craftsman god finds the visible sphere ``not at rest, but moving in an irregular and disorderly fashion,'' and ``out of disorder he brought order''---he arranges, that is, materials already there. He works by looking at a model he did not invent, ``the pattern of the unchangeable.'' And the reason given for his working at all is a moral one, in a sentence of great beauty: ``He was good, and the good can never have any jealousy of anything. And being free from jealousy, he desired that all things should be as like himself as they could be.''\endnote{Plato, \emph{Timaeus} 28a--30a, Jowett translation; wording checked 2026-07-25 at Project Gutenberg (ebook no.~1572). The same passage supplies the causal premise (``that which is created must, as we affirm, of necessity be created by a cause'') and the famous reserve (``the father and maker of all this universe is past finding out; and even if we found him, to tell of him to all men would be impossible''). Whether the making is meant temporally or as an expository device was disputed in Plato's own Academy and is disputed still; the point made here is independent of that question, since on either reading the Demiurge orders material he did not produce and consults a pattern he did not originate.} Grant the passage everything it asks and it still yields a shaper, not a source: no origin of the matter, no origin of the pattern, no account of why there is anything at all to be shaped.

Neither philosopher, in short, has any notion of a cause that supplies \emph{that things exist}. Creation from nothing is not a Greek idea; it was not rejected by the Greeks, merely never conceived. And it is worth adding that the tradition which later held it did not pretend otherwise. Aquinas, narrating the history of philosophy in the middle of his own treatment of creation, describes the ancients as advancing ``gradually, and as it were step by step'': assigning causes first to accidental change, then to substantial change, and only at the end reaching the question of being itself. ``Then others there were who arose to the consideration of `being,' as being, and who assigned a cause to things, not as `these,' or as `such,' but as `beings.'\,''\endnote{\emph{Summa theologiae} I, q.~44, a.~2, corpus; checked 2026-07-25, English at New Advent, Latin at Corpus Thomisticum. The article's conclusion---``it is necessary to say that also primary matter is created by the universal cause of things''---marks the point at which the Greek analysis is not corrected but completed: a cause of things \emph{as beings} leaves nothing outside its effect, not even the substrate the Greeks treated as underived. Aquinas's history is compressed and schematic; it is quoted as the tradition's own frank account of what its predecessors had and had not asked, not as modern historiography.} That sentence marks the seam exactly. The Greek arguments answer ``why is there motion?'' and ``why are there many so-and-sos?'' The question ``why is there anything?'' had to be posed before it could be answered, and posing it was not, in the first place, a philosophical achievement.

And yet the instruments were now on the table. A distinction between potency and act, sharp enough to show that nothing actualizes itself. A logic of ordered series, strong enough to show that simultaneous dependence must terminate in something that depends on nothing. A grammar of participation, supple enough to describe having-by-borrowing and to posit, at the source, one that is what the others merely have. What remained was to aim all three at the fact the Greeks never isolated: existence. That took a different intellectual world---one which had reasons of its own, not derived from philosophy, for suspecting that to exist is to be given.
